\begin{chapterbox}
    \chapter{Der Fall des zweitbesten Bogenschützen (WIP)}
    \label{wip-der-fall-des-zweitbesten-bogenschuetzen}
    \az{Jahr 62}
    
    \begin{center}
        Ein Neues Abenteuer von Arbon, dem Wächter der Schwarzen Archive

        Teil A der Neuen Abenteuer
    \end{center}

    Die Neuen \gend{Held*innen}{Helden} spüren eine Skralfestung im Wachsamen Wald auf. Der Schwarze Wächter Pago hat eine Rechnung mit Arbon offen.
\end{chapterbox}



Diese Fan-Geschichte ist noch in Arbeit. Ich teile sie bereits hier, damit sie für mich nicht verloren geht, falls mein Laptop spontan verschwinden sollte, und damit sie für euch nicht völlig verloren geht, falls ich spontan verschwinden sollte. Fühlt euch in letzterem Fall vollkommen frei, hierrin zu schmökern und nach Lust und Laune weiter daran zu werkeln. Aber solange ich ein aktiver Andori bleibe, lest bitte nicht zu genau nach, was ich geplant haben könnte. LG BBB :D



\newpage
\section{Prolog}

\az{Jahr -33}

Hademars Tauglichkeit zur Krahd-Hexerei wurde geprüft.



\newpage
\section{Kapitel}

\az{Jahr 62}

Die Neuen \gend{Held*innen}{Helden} saßen im Schatten von Mammutbäumen und warteten.

Sie unterhielten sich.

Fennah teilte mit, dass Morar die Skralfestung gefunden hatte.

Auf der Suche nach Unterstützung fanden sie einen Trupp von \gend{Bewahrer*innen}{Bewahrern} im Streit. Pago wollte ins Rietland, ins Dorf Geissried\sidenote{"Geiss-" angelehnt an die Ziegen auf Andor-Spielplan-Feld 40 und "-ried" angelehnt an Zwergenfürst Broms Rodung der umliegenden Wälder.}, wo Patias Schwester lebte. Ein gewisser Janis hätte Hinweise auf eine Sichtung Arbons fallen gelassen. Pagos restlicher Trupp hielt die Skralfestung für die größere Priorität.

Die Neuen \gend{Held*innen}{Helden} befreiten Chada und Thorn aus der Skralfestung und bekämpften Shron und seine Horden. Pago bekam dabei mit, wie Arbon seine Sonderfähigkeit einsetzte.

\trennzeichen

\textit{Reka reiste durch den Wachsamen Wald, auf der geistigen Spur einer Person, die sie seit Jahrzehnten für tot gehalten hatte.}


\newpage
\section{Kapitel}


Die Neuen \gend{Held*innen}{Helden} verfolgten den düsteren Drahtzieher, eine Gestalt in einem langen Umhang, umgeben von blauem Rauch und Feuer.

Pago versuchte zunächst aus dem Schatten, dann ganz offen, Arbon zu töten. Doch die übrigen Neuen \gend{Held*innen}{Helden} stellten sich auf Arbons Seite und überwältigten Pago.

Der Drahtzieher nutzte die Gelegenheit, den Spieß umzudrehen und seinerseits die Neuen \gend{Held*innen}{Helden} mit einem Schwarzen Ritter anzugreifen. Nachdem Pago ihm die Idee zur Dunklen Magie und den Drachenknochen im Grauen Gebirge gab, flüchtete Pago.

Der Drahtzieher und sein Schwarzer Ritter teilten die Neuen \gend{Held*innen}{Helden}. Der Drahtzieher verfolgte Arbon, dessen Fähigkeit er am Interessantesten fand.

Alleine hätte Arbon keine Chance gehabt, doch wurde er von einer überraschend auftauchenden Reka gerettet.

Reka und Arbon vertrieben den Drahtzieher. Sie kehrten zu den anderen Neuen \gend{Held*innen}{Helden} zurück, welche mit dem Schwarzen Ritter klargekommen waren. 

Die fünf kehrten zur Skralfestung zurück, welche soeben von den vier Grund-\gend{Held*innen}{Helden} zerstört wurde. Ende gut, alles gut?

Chada weihte Fenn ein in den Nachhall von Krahal, während Fenn Chada über Gjet na picoi grata aufklärte.

\trennzeichen

\textit{Pago drang selbst in die Schwarzen Archive ein und stahl einen uralten Holzbecher.}

\trennzeichen

\textit{Der Drahtzieher erschien am östlichen Ende des Wachsamen Waldes, wo sich eine Armee der \gend{Barbar*innen}{Barbaren} versammelt hatte. Er meldete dem Barbarenkönig, dass es an der Zeit war zu handeln. Die \gend{Bewahrer*innen}{Bewahrer} und \gend{Held*innen}{Helden} waren mit der Skralfestung beschäftigt und angeschlagen. Die Andori waren in Gefechte am Südlichen Wald verwickelt. Prinz Thorald war entführt worden, der Geist von König Brandur gebrochen. Nun musste das Land fallen.}

\textit{Der Barbarenkönig bestätigte, dass er bereits Vorhuten zum Baum der Lieder und zu verschiedenen Kapellen im Rietland ausgeschickt hatte. Die unverteidigten ketzerischen Gebäude würden fallen, sobald er das Signal gab.}


\newpage
\section{Kapitel}

Die Neuen \gend{Held*innen}{Helden} kehrten zurück zu Patias Dorf Geissried, wo Kheela ohnehin für eine Hochzeit erwartet wurde.

Kheela erfuhr von Janis und Patia, dass Pago bereits anwesend war und mit ihnen gesprochen hatte. Jetzt war er allerdings nirgendwo zu sehen.

Eine Vorhut der \gend{Barbar*innen}{Barbaren} war ebenfalls bereits anwesend und wartete auf das Signal zum Angriff. Obwohl das Signal ertönte, warteten sie das Ende der Zeremonie für den Angriff ab.

Eine Barbarin sprach Bragor in der Sprache des Lands der drei Brüder an, ob er ihr \textit{Yjotege grate} sein könnte, und war sehr enttäuscht, als er sie nicht verstand. Sprachverständnis war schließlich eines der Zeichen, an denen der große Büffel erkennbar sein sollte.

Fenn wurde von \gend{Barbar*innen}{Barbaren} als jemand aus dem Steppenland erkannt und erkannte seinerseits ihre Pläne. Obwohl er in ihrer Anwesenheit nicht frei sprechen konnte, konnte Fenn Arbon mit dem Nachhall von Krahal warnen.

Inmitten des Kampfs gegen die \gend{Barbar*innen}{Barbaren}-Vorhut stellte sich Pago einem alleine stehenden Arbon, mit Janis als Geisel. Doch Arbon war gewachsen und konnte ihn mit Varas Hilfe auch ohne die anderen überwältigen.

Pago enthüllte die drei grauen Würfel aus den Schwarzen Archiven und nutzte ihre Macht.

\trennzeichen

\textit{Es wurde sehr Meta für einen kurzen Moment.}

\trennzeichen

Ein Zeitportal öffnete sich und Griun trat samt Sphäre heraus.

Pago schnappte sich die Sphäre und verschwand im Zeitportal, gefolgt von Arbon.

Sie kämpften sich durch den Limbus, wobei Arbons Arcuballiste im Nichts verschwand.

\az{Jahr 66}

Arbon eroberte die Würfel und die Sphäre wurde beschädigt und verlor ihren Tiefsand.

Pago reiste zum Baum der Lieder, um das Hadrische Stundenglas für dessen Sand zu stehlen. Er erfuhr, dass dieses gestohlen worden war. 
\sidenote{Aus der Fan-Legende "Die verschwundene Krone", Schatzkarte 11: Hadrisches Stundenglas: \textit{Diese Gabe der Andori wurde seit der Ära des Sternenschilds sorgsam in den Schwarzen Archiven im Baum der Lieder aufbewahrt.} [...] \textit{Jemand war ins Schwarze Archiv eingedrungen und hatte ein mächtiges Artefakt gestohlen! "Zu Follas Zeiten wäre das nicht passiert", brummelte Pago mürrisch.}}

Arbon reiste zu Fenns Höhle. Der zukünftige Fenn wartete mit dem Stundenglas und weisen Worten auf ihn.

Pago erschien, füllte die Sphäre neu auf und setzte an, Arbon zu töten.

Der zukünftige Fenn tötete Pago.

Arbon überlegte, dass es dumm wäre, die Sphäre zurückzugehen.

Griun tauchte auf, bewaffnet mit der Arcuballiste, und forderte einen Tausch von Sphäre gegen Arcuballiste. Arbon bettelte vergeblich, ließ dann aber von der Sphäre ab und sich in die Vergangenheit zurückschicken.

\trennzeichen

\textit{Jarid und Trieest fanden Pagos Körper im Wald, nahe der Taverne. Trieests Lavastein leuchtete kurz auf, als Trieest neben dem Leichnam kauerte.}

\textit{Trieests Brust begann wieder zu schmerzen und ein heftiger Sturm ging durch die Bäume. Jarid bat ihn, rasch zur Taverne zu gehen.}


\begin{commentbox}
    Jarids und Trieests Geschichte geht weiter in \hyperref[2019-2020-stuermische-raetselnacht-in-der-taverne]{Stürmische Rätselnacht in der Taverne (2019 bis 2020)}
\end{commentbox}




\newpage
\section{Kapitel}


\az{Jahr 62}

Arbon erschien wieder in der Vergangenheit.

Arbon wurde vom letzten Barbar verletzt.

Die Neuen \gend{Held*innen}{Helden} klärten mit der Erzählerin ab, dass sie Arbons Geheimnis wahren würde, ehe sie auszogen, um Alarm zu schlagen. Arbon rannte dafür in Richtung der Marktbrücke.
\sidenote{Siehe den Eröffnungstext von "Der Angriff der Barbaren", wo ein verletzter Bewahrer in der Taverne Alarm schlägt.}

\trennzeichen

\textit{Die Zauberwölfe stellten sich den angreifenden \gend{Barbar*innen}{Barbaren} und sorgten so für den Sinneswandel der Wolf-Sippe.}






\newpage
\section{Epiloge}


Fenn und Bragor beschwerten sich bei Reka über ihre Lügen zu Thogger.

\az{Jahr 63}

Hademar öffnete den Zugang zur Hadrischen Unterwelt für Orweyn. 

Hademar wurde zum Dunklen Magier.

Ein Angriff auf einen Winterstein ließ ihn als Eisstatue zurück, während Orweyn davonsauste, um die Ewige Kälte auszulösen.

\az{Jahr 75}

Reka sinnierte im Grauen Gebirge vor sich hin.

Auf einmal spürte sie Hademars Präsenz.

Reka entfesselte die Grüne Hexenkatz, vergeblich.

Forn besuchte Reka. Sie suchte ängstlich in seinen Augen nach Hademars Präsenz, fühlte aber nur ihn.

Reka erzählte Forn von Hademars Stimme. Es war ihr nicht klar, ob er ihr glaubte. Er bettete sie auf ein Bett aus Sternkraut.

Reka flüsterte ihre Worte\sidenote{Siehe Forns Heldenkarte.} und Forn verließ sie wieder.

Reka und Hademar sinnierten über Arbons Meta-Macht.

Chada besuchte Reka.

Reka starb.
\sidenote{Siehe Chadas Heldenkarte.}








\begin{commentbox}
    Weiter geht es in \hyperref[wip-kampf-in-den-flusslanden]{Kampf in den Flusslanden (WIP)}.
\end{commentbox}
